\section{Introduction}\label{sec:intro}
Checking a code change by running the old and the new version on the same inputs is a long-standing practice~\cite{evans2007,palikareva2016,jin2010automated}, and recent work uses large language models to generate the inputs~\cite{etemadi2024mokav,pradel2025testora,hu2026difftestgen,liu2025evaluating}. A difference found this way is actionable. The absence of a difference is not: it holds only for the executions that were observed, and a change that was never reached by any input produces the same silence as a change that was exercised and found harmless.

This paper describes V7-B, a research prototype built around one design rule: a ``no difference'' outcome must carry a record of how much of the change was exercised, and must be withheld when that record is inadequate or when nondeterminism, resource limits or engine gaps make the comparison unreliable. The prototype is implemented for Python functions and evaluated only on development data. We state the scope early because the evidence is limited: all 321 evaluation pairs are exposed or non-pristine, no independent holdout exists, execution isolation was exercised on macOS only, and the central design element (the sufficiency gate) lowers false assurance only by abstaining, without improving detection (Section~\ref{sec:results}).

The contribution is therefore an explicit design, an implementation status report, and an exploratory Phase 2 result set that includes negative and null findings, offered so that a confirmatory study can be designed on firmer ground. Section~\ref{sec:rqs} lists the claims we do and do not make.
